\section[Cote 42 : familles cohérentes de germes, le lemme 1 et le lemme 2]{Cote 42 : familles cohérentes de germes, le lemme 1 et le lemme 2\protect\verifie{Folder42}}\label{sec:42}

La cote 42, \emph{Réseaux d'anneaux locaux : notes manuscrites (s.d.)}, six pages que l'inventaire date \q{[années 1960-1970]}, demande si une application entre préschémas localement noethériens est déterminée, et peut être construite, point par point, à partir des anneaux locaux reliés par les générisations. La page~3 ramène la question à une proposition sur les faisceaux : sous une condition~(H) d'injectivité des flèches de générisation entre fibres, les sections globales sont les familles cohérentes de germes. Le lemme~1 en porte le contenu, et la page le démontre en une ligne. Le lieu considéré est stable par générisation, et \q{il suffit de montrer qu'il est \emph{constructible}, donc que s'il contient $x$, il contient les pts d'un ouvert non vide de $\bar x$} (\lot{42}{1}{3}). Pour le faisceau structural, la page~4 enchaîne : \q{Il reste à vérifier que les conditions de la prop.\ sont vérifiées}. Elle isole un fait d'algèbre commutative, le lemme~2, dont la démonstration occupe le bas de la page~4 et le haut de la page~5 (\lot{42}{1}{4}). La page glose le noyau de $A \to A_{\mathfrak p}$ comme \q{l'annulateur de $A - p$}, le premier mot étant lu avec doute. Elle traite la seconde assertion en se ramenant au cas $f = 1$, \q{car $A_p = (A_f)_{p A_f}$} (\lot{42}{1}{5}).

\begin{enonce}[Proposition et lemme 1, page 3]
Soient $X$ un espace topologique noethérien et sobre, et $F$ un faisceau d'ensembles sur $X$ vérifiant
\begin{itemize}
\item[(H)] pour tout $x \in X$, il existe un ouvert $W \ni x$ tel que $\rho_{yx} : F_y \to F_x$ soit injective pour tout $y \in W \cap \overline{\{x\}}$.
\end{itemize}
\begin{enumerate}
\item (Lemme~1) Pour toute famille cohérente $(\xi_x)$ et toute section $s \in F(U)$ sur un ouvert $U$, l'ensemble $E = \{x \in U : s_x = \xi_x\}$ est ouvert.
\item (Proposition) L'application $\Gamma(X, F) \to \varprojlim_x F_x$, qui envoie une section sur la famille de ses germes, est bijective.
\end{enumerate}
\end{enonce}

\begin{enonce}[Lemme 2, pages 4 et 5]
Soient $A$ un anneau noethérien et $\mathfrak p$ un idéal premier de $A$.
\begin{enumerate}
\item Il existe $f \in A \setminus \mathfrak p$ tel que $A_f \to A_{\mathfrak p}$ soit injectif.
\item Si de plus $A \to A_{\mathfrak p}$ est injectif, alors $A_{\mathfrak q} \to A_{\mathfrak p}$ est injectif pour tout idéal premier $\mathfrak q \supseteq \mathfrak p$.
\item (Forme de la page) Il existe $f \in A \setminus \mathfrak p$ tel que, pour tout idéal premier $\mathfrak Q$ de $A_f$ dont la trace sur $A$ contient $\mathfrak p$, $(A_f)_{\mathfrak Q} \to A_{\mathfrak p}$ soit injectif.
\end{enumerate}
\end{enonce}

\begin{enonce}[Lemme 1 et corollaire 1, pages 3 et 4, cas affine]
Soit $A$ un anneau noethérien. Appelons famille cohérente une famille $(\xi_{\mathfrak p})$, $\xi_{\mathfrak p} \in A_{\mathfrak p}$, indexée par les idéaux premiers, telle que $\xi_{\mathfrak q} \mapsto \xi_{\mathfrak p}$ par $A_{\mathfrak q} \to A_{\mathfrak p}$ chaque fois que $\mathfrak p \subseteq \mathfrak q$.
\begin{enumerate}
\item (Lemme~1) Pour toute famille cohérente $(\xi_{\mathfrak p})$ et tous $a, s \in A$, l'ensemble des $\mathfrak p$ tels que $s \notin \mathfrak p$ et $\xi_{\mathfrak p} = a/s$ est ouvert dans $\operatorname{Spec} A$.
\item (Corollaire~1) L'application $a \mapsto (a/1)_{\mathfrak p}$ est une bijection de $A = \Gamma(\operatorname{Spec} A, \mathcal O)$ sur l'ensemble des familles cohérentes.
\end{enumerate}
\end{enonce}

\noindent\emph{Conventions.} Pour $x, y \in X$, on écrit $x \rightsquigarrow y$ lorsque $y \in \overline{\{x\}}$ : $y$ est une spécialisation de $x$, $x$ une générisation de $y$. Tout ouvert contenant $y$ contient alors $x$, d'où une application $\rho_{yx} : F_y \to F_x$, compatible aux germes : $\rho_{yx}(s_y) = s_x$ pour toute section $s$ définie au voisinage de $y$. Une famille $(\xi_x)$, $\xi_x \in F_x$, est cohérente si $\rho_{yx}(\xi_y) = \xi_x$ chaque fois que $x \rightsquigarrow y$. Dans $\operatorname{Spec} A$, $\mathfrak p \rightsquigarrow \mathfrak q$ équivaut à $\mathfrak p \subseteq \mathfrak q$. Un ouvert non vide de l'irréductible $\overline{\{x\}}$ contient son point générique $x$ ; la condition~(H) dit donc la même chose que la page, \q{un ouvert non vide} de $\overline{\{x\}}$ sur lequel les $F_y \to F_x$ sont injectives. Pour une partie multiplicative $S$ de $A$, on note $a \mapsto a/1$ le morphisme $A \to S^{-1}A$ ; on rappelle que $a/1 = 0$ si et seulement s'il existe $s \in S$ avec $sa = 0$, que tout $z \in S^{-1}A$ vérifie $z \cdot (s/1) = a/1$ pour un $a \in A$ et un $s \in S$, et que si $g : A \to B$ envoie $S$ dans les inversibles de $B$, il existe un unique $\bar g : S^{-1}A \to B$ avec $\bar g(a/1) = g(a)$. Les flèches $A_f \to A_{\mathfrak p}$ et $A_{\mathfrak q} \to A_{\mathfrak p}$ sont de cette forme, $g$ étant le morphisme canonique $A \to A_{\mathfrak p}$.

\begin{lemme}\label{lem:42-ouvert}
Soient $X$ un espace noethérien et sobre et $E \subseteq X$ une partie stable par générisation telle que, pour tout $x \in E$, il existe un ouvert $W \ni x$ avec $W \cap \overline{\{x\}} \subseteq E$. Alors $E$ est ouvert.
\end{lemme}

\begin{proof}
Montrons d'abord, par récurrence noethérienne sur les fermés, que tout fermé $Z$ tel que $Z \subseteq \overline{Z \setminus E}$ ne rencontre pas $E$. Soit $Z$ un tel fermé, et supposons l'assertion vraie pour les fermés strictement contenus dans $Z$. Si $Z$ est irréductible, soit $z$ son point générique. Si $Z$ rencontrait $E$ en un point $y$, on aurait $z \rightsquigarrow y$, donc $z \in E$ par stabilité ; l'hypothèse donne un ouvert $W \ni z$ avec $W \cap Z \subseteq E$. Mais $z \in Z \subseteq \overline{Z \setminus E}$, donc $W$ rencontre $Z \setminus E$, en un point de $W \cap Z$ qui n'est pas dans $E$ : contradiction. Si $Z$ n'est pas irréductible (et n'est pas vide), il existe deux fermés $z_1, z_2$ avec $Z \subseteq z_1 \cup z_2$ et $Z \not\subseteq z_i$. Posons $Z_i = \overline{Z \cap z_i \setminus E}$ ; c'est un fermé contenu dans $Z \cap z_i$, donc strictement contenu dans $Z$, et $Z_i \subseteq \overline{Z_i \setminus E}$ puisque $Z \cap z_i \setminus E \subseteq Z_i \setminus E$. Par hypothèse de récurrence, $Z_i \cap E = \varnothing$. Or $Z \setminus E \subseteq (Z \cap z_1 \setminus E) \cup (Z \cap z_2 \setminus E)$, d'où $Z \subseteq \overline{Z \setminus E} \subseteq Z_1 \cup Z_2$, et $Z \cap E = \varnothing$.

Appliquons cela à $T = \overline{X \setminus E}$, qui vérifie $T \setminus E = X \setminus E$, donc $T \subseteq \overline{T \setminus E}$ : il vient $T \cap E = \varnothing$, c'est-à-dire que $X \setminus E$ est fermé.
\end{proof}

\begin{proof}[Démonstration du lemme 1 de la page 3]
\emph{$E$ est stable par générisation.} Soient $x \in E$ et $y \rightsquigarrow x$. Comme $U$ est ouvert et contient $x$, il contient $y$, et $s_y = \rho_{xy}(s_x) = \rho_{xy}(\xi_x) = \xi_y$ par cohérence ; donc $y \in E$.

\emph{Condition locale.} Soit $x \in E$, et soit $W$ l'ouvert que fournit~(H) en $x$. L'ouvert $W \cap U$ contient $x$. Pour $y \in W \cap U \cap \overline{\{x\}}$, on a $x \rightsquigarrow y$, et $\rho_{yx}(s_y) = s_x = \xi_x = \rho_{yx}(\xi_y)$ ; l'injectivité de $\rho_{yx}$ donne $s_y = \xi_y$, donc $y \in E$. Le lemme~\ref{lem:42-ouvert} conclut.
\end{proof}

\begin{proof}[Démonstration de la proposition]
L'injectivité est l'axiome de séparation des faisceaux : deux sections globales qui ont les mêmes germes en tout point sont égales. Soit $(\xi_x)$ une famille cohérente. Pour chaque $x$, le germe $\xi_x$ provient d'une section $s^x$ sur un ouvert $U_x \ni x$ ; par le lemme~1, $V_x = \{y \in U_x : s^x_y = \xi_y\}$ est un ouvert, qui contient $x$. Les restrictions $s^x|_{V_x}$ ont, sur $V_x \cap V_{x'}$, les mêmes germes $\xi_y$ en tout point ; elles y coïncident donc, et se recollent, les $V_x$ recouvrant $X$, en une section globale $s$. Pour tout $x$, $s_x = (s^x|_{V_x})_x = \xi_x$.
\end{proof}

\begin{lemme}\label{lem:42-lift}
Soient $S \subseteq A$ une partie multiplicative et $g : A \to B$ un morphisme d'anneaux envoyant $S$ dans les inversibles de $B$. Si, pour tout $a \in A$, $g(a) = 0$ entraîne $a/1 = 0$ dans $S^{-1}A$, alors $\bar g : S^{-1}A \to B$ est injectif.
\end{lemme}

\begin{proof}
Comme $\bar g$ est additif, il suffit de montrer que son noyau est nul. Soit $z \in S^{-1}A$ avec $\bar g(z) = 0$, et écrivons $z \cdot (s/1) = a/1$ avec $a \in A$ et $s \in S$. En appliquant $\bar g$, on obtient $g(a) = \bar g(z)\, g(s) = 0$, donc $a/1 = 0$ par hypothèse, c'est-à-dire $z \cdot (s/1) = 0$. Comme $s/1$ est inversible dans $S^{-1}A$, $z = 0$.
\end{proof}

\begin{lemme}\label{lem:42-tue}
Soient $A$ un anneau noethérien et $\mathfrak p$ un idéal premier. Il existe $f \notin \mathfrak p$ tel que $fy = 0$ pour tout $y$ du noyau $J$ de $A \to A_{\mathfrak p}$.
\end{lemme}

\begin{proof}
Le noyau $J$ est l'ensemble des $a \in A$ tels que $sa = 0$ pour \emph{au moins un} $s \notin \mathfrak p$. L'anneau $A$ étant noethérien, l'idéal $J$ est engendré par un nombre fini d'éléments $x_1, \dots, x_n$. Pour chaque $i$, l'appartenance $x_i \in J$ fournit un $s_i \notin \mathfrak p$ tel que $s_i x_i = 0$. Posons $f = s_1 \cdots s_n$ ; comme $\mathfrak p$ est premier, son complémentaire est stable par produit, et $f \notin \mathfrak p$. L'ensemble des $y \in A$ tels que $fy = 0$ est un idéal de $A$ ; il contient chaque $x_i$, car $f = s_i c_i$ avec $c_i = \prod_{j \neq i} s_j$, donc $f x_i = c_i (s_i x_i) = 0$ ; il contient donc l'idéal engendré par les $x_i$, qui est $J$.
\end{proof}

\begin{proof}[Démonstration du lemme 2]
(1) Soit $f$ donné par le lemme~\ref{lem:42-tue}. Appliquons le lemme~\ref{lem:42-lift} à $S = \{f^n\}$ et au morphisme canonique $g : A \to A_{\mathfrak p}$. Si $g(a) = 0$, alors $a \in J$, donc $fa = 0$ ; comme $f \in S$, cela signifie $a/1 = 0$ dans $A_f$. Le morphisme $A_f \to A_{\mathfrak p}$ est donc injectif.

(2) Appliquons le lemme~\ref{lem:42-lift} à $S = A \setminus \mathfrak q$ et au morphisme canonique $g : A \to A_{\mathfrak p}$, qui inverse $S$ puisque $\mathfrak p \subseteq \mathfrak q$. Si $g(a) = 0$, alors $a = 0$ par l'injectivité supposée de $g$, et a fortiori $a/1 = 0$ dans $A_{\mathfrak q}$.

(3) Prenons $f$ comme en~(1), et soit $\mathfrak Q$ un idéal premier de $A_f$ dont la trace $\mathfrak q$ sur $A$ contient $\mathfrak p$. Le morphisme $h : A_f \to A_{\mathfrak p}$ envoie $A_f \setminus \mathfrak Q$ dans les inversibles : écrivons $y \in A_f \setminus \mathfrak Q$ sous la forme $y \cdot (f^n/1) = a/1$ ; si $a$ était dans $\mathfrak p$, donc dans $\mathfrak q$, $a/1 = y \cdot (f^n/1)$ serait dans $\mathfrak Q$, ce qui est exclu puisque $y \notin \mathfrak Q$ et que $f^n/1$, inversible, n'est pas dans $\mathfrak Q$. Donc $a \notin \mathfrak p$, $h(y)\,h(f^n/1) = a/1$ est inversible dans $A_{\mathfrak p}$, et $h(y)$ l'est aussi. Le morphisme $(A_f)_{\mathfrak Q} \to A_{\mathfrak p}$ est donc défini, et il est injectif par le lemme~\ref{lem:42-lift} appliqué à $h$, qui est injectif par~(1) : c'est l'argument de~(2) pour l'anneau $A_f$.
\end{proof}

\begin{proof}[Démonstration du lemme 1 et du corollaire 1, cas affine]
Pour le lemme~1, soit $E$ l'ensemble des $\mathfrak p$ tels que $s \notin \mathfrak p$ et $\xi_{\mathfrak p} \cdot (s/1) = a/1$. Il est stable par générisation : si $\mathfrak q \in E$ et $\mathfrak p \subseteq \mathfrak q$, alors $s \notin \mathfrak p$, et l'image par $A_{\mathfrak q} \to A_{\mathfrak p}$ de l'égalité $\xi_{\mathfrak q} \cdot (s/1) = a/1$ est $\xi_{\mathfrak p} \cdot (s/1) = a/1$. Soit $\mathfrak p \in E$, et $f$ donné par le lemme~\ref{lem:42-tue}. L'ouvert $D(f) \cap D(s)$ contient $\mathfrak p$. Soit $\mathfrak q$ un premier de cet ouvert qui contient $\mathfrak p$ ; le lemme~\ref{lem:42-lift}, appliqué à $S = A \setminus \mathfrak q$, montre que $A_{\mathfrak q} \to A_{\mathfrak p}$ est injectif : si $a'/1 = 0$ dans $A_{\mathfrak p}$, alors $fa' = 0$ avec $f \notin \mathfrak q$, donc $a'/1 = 0$ dans $A_{\mathfrak q}$. Or $\xi_{\mathfrak q} \cdot (s/1) - a/1$ a pour image $\xi_{\mathfrak p} \cdot (s/1) - a/1 = 0$ ; il est donc nul, et $\mathfrak q \in E$. Le lemme~\ref{lem:42-ouvert} s'applique, $\operatorname{Spec} A$ étant noethérien et sobre.

Pour le corollaire~1, on utilise la description du faisceau structural comme faisceau des familles $(\sigma_{\mathfrak p}) \in \prod_{\mathfrak p \in U} A_{\mathfrak p}$ qui sont localement des fractions --- pour tout $\mathfrak p \in U$, il existe un voisinage $V$ de $\mathfrak p$ et $a, s \in A$ avec $s \notin \mathfrak q$ et $\sigma_{\mathfrak q} = a/s$ pour tout $\mathfrak q \in V$ --- et le fait que $a \mapsto (a/1)_{\mathfrak p}$ est une bijection de $A$ sur les sections globales de ce faisceau. L'injectivité de l'application du corollaire en découle. Soit $(\xi_{\mathfrak p})$ une famille cohérente. Pour chaque $\mathfrak p$, écrivons $\xi_{\mathfrak p} \cdot (s/1) = a/1$ avec $s \notin \mathfrak p$ ; l'ensemble $E$ du lemme~1 pour ces $a$ et $s$ est un voisinage ouvert de $\mathfrak p$, sur lequel $\xi_{\mathfrak q} = a/s$. La famille $(\xi_{\mathfrak p})$ est donc localement une fraction, c'est une section globale du faisceau structural, et elle provient d'un $a \in A$.
\end{proof}

\begin{remarque}
Le lemme~\ref{lem:42-ouvert} est le \q{donc} de la page~3. De la constructibilité, la page passe aussitôt à la propriété \q{s'il contient $x$, il contient les points d'un ouvert non vide de $\bar x$}. Le lemme prend pour hypothèse cette propriété, jointe à la stabilité par générisation, et la constructibilité n'intervient pas dans la démonstration. La bibliothèque utilisée contient l'énoncé \q{constructible et stable par générisation, donc ouvert} pour les spectres d'anneaux. Elle ne contient pas le critère de constructibilité des espaces noethériens par lequel la page y arrive, et le lemme fait les deux pas à la fois.
\end{remarque}

\begin{remarque}
Les énoncés sont démontrés pour un espace noethérien, et le corollaire~1 pour $X = \operatorname{Spec} A$ affine. C'est le cas auquel la page se ramène (\q{La question est locale. On peut donc supposer $X$ noethérien}). La page~3 écrit \q{un voisinage $U$ de $\bar x$ tel que pour tt $y \in \bar x \cap U$}. Un ouvert contenant tout $\overline{\{x\}}$ donnerait une condition plus forte que~(H), qui l'entraîne. La démonstration de la page, comme celle-ci, n'emploie que la forme faible, \q{un ouvert non vide $U$ de $\bar x$}. L'injectivité de $\Gamma(X,F) \to \varprojlim F_x$ ne demande ni~(H) ni l'hypothèse noethérienne, et celle de $A \to \varprojlim A_{\mathfrak p}$ vaut pour tout anneau commutatif.
\end{remarque}

\begin{remarque}
L'hypothèse noethérienne n'intervient pas dans~(2) du lemme~2, qui vaut pour tout anneau commutatif. Dans~(1), elle sert seulement à rendre de type fini le noyau $J$ de $A \to A_{\mathfrak p}$. Ce noyau est formé des éléments annulés par \emph{un} élément de $A \setminus \mathfrak p$. L'annulateur de $A \setminus \mathfrak p$, que glose la page, est en général plus petit et ne convient pas. La démonstration de la page prend un annulateur par générateur, et suit donc la bonne description.
\end{remarque}

\begin{remarque}
Dans la forme~(3) du lemme~2, l'hypothèse que la trace de $\mathfrak Q$ contienne $\mathfrak p$ est celle de la page, \q{provenant d'un idéal $q' \supset p$ ne rencontrant pas $f$} ; elle rend la flèche $(A_f)_{\mathfrak Q} \to A_{\mathfrak p}$ définie. La démonstration n'a pas besoin de l'identification $A_{\mathfrak p} = (A_f)_{\mathfrak p A_f}$ par laquelle la page se ramène au cas $f = 1$. Le corollaire~1 n'utilise du lemme~2 que le lemme~\ref{lem:42-tue}. Un même $f$ tue le noyau de $A \to A_{\mathfrak p}$ et rend injectives toutes les flèches $A_{\mathfrak q} \to A_{\mathfrak p}$ pour $\mathfrak q \in D(f) \cap V(\mathfrak p)$, ce qui est~(H) pour le faisceau structural. Dans la démonstration du lemme~1, l'argument donne l'égalité $s_y = \xi_y$, là où la page écrit \q{d'où $s_x = \xi_y$}.
\end{remarque}
